\documentclass[11pt]{article}
\usepackage[T1]{fontenc}
\usepackage{lmodern,amsmath,amssymb,amsthm,mathtools}
\usepackage[margin=1in]{geometry}
\usepackage{microtype}
\usepackage[hidelinks]{hyperref}
\newtheorem{theorem}{Theorem}[section]
\newtheorem{lemma}[theorem]{Lemma}
\newtheorem{proposition}[theorem]{Proposition}
\newtheorem{corollary}[theorem]{Corollary}
\newcommand{\R}{\mathbb R}\newcommand{\E}{\mathbb E}
\newcommand{\Hh}{\mathcal H}\newcommand{\U}{\mathcal U}
\newcommand{\code}[1]{{\small\nolinkurl{#1}}}
\title{Gaussian prediction and recovery:\\Hermite rank, observation dimension and temporal obstructions}
\author{Project Stardust}
\date{Reviewed formalization companion, October 29, 2026}
\begin{document}
\maketitle
\begin{abstract}
For a square-integrable function of one standard Gaussian coordinate,
conditioning on any Gaussian observation family transports its Hermite
expansion to the orthogonal projection of that coordinate. We reconstruct
the full $L^2$ prediction formula, exact prediction-energy series and a
quantitative leading-rank remainder. Applied to the actual full past of
fractional Brownian motion, these identities give explicit lower bounds
with exponent $2m\min(H,1-H)$ for any nonzero Hermite coefficient of rank
$m$, and exclude higher-power increment bounds at zero information.
The inverse-number potential divides positive-rank coefficients by their
rank, so it preserves this obstruction despite being compact on a single
Gaussian observation space. A concrete quarter-Hurst covariance witness
also separates full-past prediction from current-value prediction. An
exact multi-target quartic calculation gives independent targets whose
shared-observation residuals have negative cross-covariance, so block error
cannot be assembled from diagonal residual energies alone. Exact finite-block
retained variances include the total-variance correction terms; shifted
scalar quartic variance has a strict minimum at zero shift.
The projection and Hermite mechanisms are classical, and the two
fractional covariance regimes are already known in primary literature.
This paper supplies complete proofs and exact frozen formal
correspondences; no priority claim is made.
The same method controls terminal-aware Brownian bracket errors for bounded Borel weights. We also reconstruct the full joint source/sample recovery identity, compact recovery exactly at finite observed source dimension, and a positive-chaos compactness criterion. Infinite observed dimension has distance one from compact recovery approximants; even one-source compact recovery has divergent squared image energy.
An accelerated Mehler limit retains full-past Markov kernels while
losing initial continuity and a dense right norm-generator domain.
\end{abstract}

\section{Setting and main prediction formula}
Let $X:\Hh\to L^2(\Omega,P)$ be an isonormal real Gaussian family on a
real separable Hilbert space: $X$ is linear, centered, jointly Gaussian,
and $\E[X(h)X(k)]=\langle h,k\rangle$. Fix $u\in\Hh$ with $\|u\|=1$.
Let $(h_i)_{i\in I}$ be any observation family, and put
\[
 M=\overline{\operatorname{span}}\{h_i:i\in I\},\quad
 p=P_Mu,\quad v=\|p\|^2\in[0,1],\quad
 \mathcal A=\sigma(X(h_i):i\in I).
\]
We use $L^2$ equivalence classes; observation sigma-algebras are not
identified literally when they differ on null sets.

Define Wick polynomials $W_n(h)$ by the generating identity
\begin{equation}\label{eq:gen}
 \exp(tX(h)-t^2\|h\|^2/2)
 =\sum_{n=0}^\infty\frac{t^n}{n!}W_n(h).
\end{equation}
Thus $W_0(h)=1$, $W_1(h)=X(h)$ and $W_2(h)=X(h)^2-\|h\|^2$.
At unit $u$, these are the probabilists' Hermite polynomials
$\operatorname{He}_n(X(u))$.

\begin{theorem}\label{thm:energy}
Every $f\in L^2(\sigma(X(u)),P)$ has a unique expansion
\begin{equation}\label{eq:terminal}
 f=\sum_{n\geq0}a_n W_n(u),\qquad
 a_n=\frac{\E[fW_n(u)]}{n!},\qquad
 \|f\|_2^2=\sum_{n\geq0}a_n^2n!.
\end{equation}
Conditioning gives the norm-convergent series and exact energy
\begin{equation}\label{eq:prediction}
 \E[f\mid\mathcal A]=\sum_{n\geq0}a_n W_n(p),\qquad
 \|\E[f\mid\mathcal A]\|_2^2=\sum_{n\geq0}a_n^2n!v^n.
\end{equation}
If $a_n=0$ for $n<m$, and $A_m=a_m^2m!$, then
\begin{align}
 A_m v^m&\leq\|\E[f\mid\mathcal A]\|_2^2
                    \leq v^m\|f\|_2^2,\label{eq:rankbounds}\\
 0&\leq\|\E[f\mid\mathcal A]\|_2^2-A_mv^m
                    \leq v^{m+1}\|f\|_2^2.\label{eq:remainder}
\end{align}
For any varying observation families with $v>0$ and $v\to0$,
\begin{equation}\label{eq:ranklimit}
 \frac{\|\E[f\mid\mathcal A]\|_2^2}{v^m}\longrightarrow A_m,
 \qquad
 0\leq\frac{\|\E[f\mid\mathcal A]\|_2^2}{v^m}-A_m
       \leq v\|f\|_2^2.
\end{equation}
\end{theorem}

The exact nonlinear prediction mechanism is established Gaussian
theory. McElroy and Das~\cite{MD} give Hermite-based nonlinear prediction
formulas for Gaussian transformations with finite data. Here the
Hilbert projection describes arbitrary observation families and makes
the $L^2$ completion and rank remainder explicit. This is a
formalization exposition, not a claim to have discovered nonlinear
Gaussian prediction.

\section{Completeness and the conditional Wick identity}
\begin{lemma}\label{lem:hermite}
For all $h,k\in\Hh$ and nonnegative integers $n,m$,
\[
 \E[W_n(h)W_m(k)]=\mathbf1_{n=m}\,n!\langle h,k\rangle^n.
\]
The polynomials $W_n(u)$ have dense linear span in
$L^2(\sigma(X(u)),P)$.
\end{lemma}
\begin{proof}
The Gaussian exponential formula gives
\[
 \E\left[e^{sX(h)-s^2\|h\|^2/2}
         e^{tX(k)-t^2\|k\|^2/2}\right]=e^{st\langle h,k\rangle}.
\]
Taking coefficients yields orthogonality and its exact constant.
Coefficient extraction under expectation is valid: on any compact set
of real or complex parameters, a Gaussian exponential with a fixed
larger real linear coefficient dominates each differentiated expression
in $L^1$, by the Gaussian exponential moment formula.

To prove completeness rather than assume it, identify the terminal
space with $L^2(\gamma)$, $\gamma$ standard Gaussian measure. If $g$ is
orthogonal to all polynomials, consider
$F(z)=\int g(x)e^{zx}\,d\gamma(x)$. Cauchy--Schwarz and Gaussian
exponential moments make this function entire, with derivatives obtained
under the integral uniformly on compact parameter sets. All its Taylor
coefficients vanish, so $F$ is identically zero. In particular the finite
signed measure $g\gamma$ has zero Fourier transform.
For completeness, Fourier uniqueness here follows by Gaussian
convolution: Fourier inversion for a Gaussian and Fubini show each
Gaussian convolution of this measure is zero; taking an approximate
identity gives zero integrals against every compactly supported
continuous function, and hence the measure is zero. Thus $g=0$.
The Wick polynomials span all polynomials, because their leading
coefficients are one. The orthogonal complement is therefore zero,
proving density and, by Hilbert-space orthogonal expansion,
\eqref{eq:terminal}.
\end{proof}

\begin{lemma}\label{lem:conditional}
$\E[W_n(u)\mid\mathcal A]=W_n(p)$ for every $n$.
\end{lemma}
\begin{proof}
Since $u=p+(u-p)$ and $u-p\perp M$, $X(u-p)$ is independent of every
finite vector of observations: their joint Gaussian characteristic
function factors because all cross-covariances are zero. A pi-lambda
argument extends independence to $\mathcal A$. Also $X(p)$ is
$\mathcal A$-measurable modulo null sets. Indeed $p$ is a norm limit of
finite linear combinations of observation directions, so $X(p)$ is an
$L^2$ limit of measurable variables and has a measurable representative.
Consequently
\[
 \E[e^{tX(u)-t^2/2}\mid\mathcal A]
 =e^{tX(p)-t^2\|p\|^2/2}.
\]
The residual exponential has expectation one, giving the equality.
The exponential series is analytic in $L^2$ on compact parameter sets:
the same Gaussian domination as in Lemma~\ref{lem:hermite}, now using
squared norms, justifies differentiating. Conditional expectation is a
bounded linear map on $L^2$, so differentiation commutes with it.
Taking the $n$th derivative at zero proves the claim.
\end{proof}

\begin{proof}[Proof of Theorem~\ref{thm:energy}]
Lemma~\ref{lem:hermite} gives the terminal expansion and Parseval.
Conditional expectation is an $L^2$ contraction, so it commutes with
its partial-sum limit. Lemma~\ref{lem:conditional} yields the first
formula in \eqref{eq:prediction}. The terms are mutually orthogonal
and have squared norms $a_n^2n!v^n$; taking norms of finite sums and
passing to the limit gives the energy formula.
If the coefficients below $m$ vanish, all terms start at $m$. Since
$0\leq v\leq1$, their sum is at most $v^m\sum_n a_n^2n!$, and at
least its $m$th term. After subtracting that term, every remaining
power is at most $v^{m+1}$, proving \eqref{eq:remainder}. For $v>0$
division by $v^m$ proves the last bound in \eqref{eq:ranklimit}; the
squeeze theorem proves its limit.
\end{proof}

If $v=0$, the predictor is $a_0=\E f$, with $0^0=1$ in the series.
No normalized-rank quotient is then used. If $f$ is centered and
nonzero, Parseval ensures a least positive $m$ with $a_m\ne0$; this
is its Hermite rank. The upper bound in \eqref{eq:rankbounds} is
attained for a pure rank-$m$ polynomial. No assumption that $f$ itself
is polynomial is needed anywhere.

\section{Actual fractional pasts and the two covariance regimes}
Let $B^H$ be a continuous fractional Brownian motion with $0<H<1$,
normalized by
\[
 R_H(s,t)=\E[B_s^HB_t^H]
 =\tfrac12(s^{2H}+t^{2H}-|t-s|^{2H}),\qquad s,t\geq0.
\]
It has the Gaussian Hilbert realization $B_t^H=X(u_t)$,
$\langle u_s,u_t\rangle=R_H(s,t)$. In particular $\|u_1\|=1$.
For $0\leq t\leq1$ put
\[
 \mathcal F_t^H=\sigma(B_s^H:0\leq s\leq t),\quad
 P_t=\E[\,\cdot\mid\mathcal F_t^H],\quad
 v_H(t)=\|P_tB_1^H\|_2^2,\quad \kappa_H=2\min(H,1-H).
\]
These are conditional expectations onto the entire actual observed
past. In the frozen construction a continuous modification realizes
the Hilbert coordinates. Coordinatewise almost-sure equality suffices
to identify their $L^2$ observation spaces: the sets agreeing modulo
null sets form a sigma-algebra containing each coordinate preimage.
Closure under countable unions only uses countable unions of null
sets. Thus the augmented observation sigma-algebras agree, and the
conditional expectations agree as $L^2$ classes. No common exceptional
set for an uncountable family or literal equality of raw sigma-algebras
is assumed.

\begin{proposition}\label{prop:variance}
For $0<t\leq1$,
\begin{equation}\label{eq:variance}
 v_H(t)\geq\tfrac14t^{\kappa_H}.
\end{equation}
Consequently, for $f\in L^2(\sigma(B_1^H),P)$ with terminal Hermite
coefficients $a_n$ and every nonnegative integer $m$,
\begin{equation}\label{eq:fractional}
 \|P_tf\|_2^2\geq\frac{a_m^2m!}{4^m}t^{\kappa_Hm}.
\end{equation}
If $f$ is centered, $a_m\ne0$, and $\alpha>\kappa_Hm$, then for every
real $C$ and every $\delta>0$ some $0<t<\min(\delta,1)$ satisfies
\begin{equation}\label{eq:nohigher}
 C t^\alpha<\|P_tf-P_0f\|_2^2.
\end{equation}
\end{proposition}
\begin{proof}
The current observation is one member of the past. Orthogonality of
conditional projection therefore gives
\[
 \langle P_tB_1^H,B_t^H\rangle=R_H(1,t),\qquad
 v_H(t)\geq R_H(1,t)^2/t^{2H}.
\]
Write $r=R_H(1,t)=(1+t^{2H}-(1-t)^{2H})/2$.
Since $(1-t)^{2H}\leq1$, $r\geq t^{2H}/2$.
Also $r\geq t/2$: if $2H\leq1$ use $t^{2H}\geq t$ in the preceding
bound, while if $2H\geq1$ use $(1-t)^{2H}\leq1-t$ in the exact formula.
The two bounds give
\[
 v_H(t)\geq\tfrac14\max(t^{2H},t^{2-2H})
          =\tfrac14t^{2\min(H,1-H)}.
\]
Theorem~\ref{thm:energy} and the single nonnegative rank-$m$ energy term
now prove \eqref{eq:fractional}.

The zero-time past is generated by $B_0^H=0$, so for centered $f$,
$P_0f=0$. Put $K=a_m^2m!/4^m>0$ and $\beta=\kappa_Hm$.
Since $\alpha>\beta$, $Ct^{\alpha-\beta}\to0$ as $t\downarrow0$.
Choose $t$ smaller than $\delta$ and one, with this expression less
than $K$. Multiplying by $t^\beta$ and using
\eqref{eq:fractional} gives \eqref{eq:nohigher}.
\end{proof}

The lower constant $1/4$ is not claimed optimal, and this proof does not
give an asymptotic equivalence for $v_H$. Sottinen and Viitasaari's
primary prediction-law paper~\cite{SV} gives explicit full-past
conditional laws and no-information covariance asymptotics in the
two regimes. Our elementary lower bound uses only the covariance and
one current observation to obstruct any faster full-past decay. It
does not claim discovery of those exponents or replace the explicit
prediction kernel.

\subsection*{An initial projected-leading-term remainder diagnostic}
The variance lower bound also tests a specific proposed remainder,
without assuming a general fractional It\^o or Taylor theorem.
Let $u_t$ be the actual fractional Gaussian directions, so
$B_t^H=X(u_t)$ and $u_0=0$. Write
\[
 K_t=\overline{\operatorname{span}}\{u_s:0\leq s\leq t\},
 \qquad Q_t:\mathcal H\to K_t
\]
for the source projection, distinct from conditional expectation
on random variables. Let $F_t=\E[B_1^H\mid\mathcal F_t^H]$ be the
actual full-past forecast already studied. Define the proposed
initial scalar leading term and its remainder by
\[
 L_0(t)=\langle Q_0u_1,u_t-u_0\rangle,\qquad
 \mathcal R_t=F_t-F_0-L_0(t)\mathbf1.
\]
Because the only observed direction at zero is $u_0=0$, one has
$K_0=\{0\}$ and $Q_0=0$. Also $F_0=0$ by centering and the
trivial initial Gaussian observation. Hence $L_0(t)=0$ and
$\mathcal R_t=F_t$ exactly.

For every $0<H<1$, put $\alpha_H=\min(4H,2H+1)$.
Both candidates exceed $2H$, so $\alpha_H>2H$. The proved
current-coordinate lower bound gives, for $0<t\leq1$,
\[
 \|\mathcal R_t\|_2^2\geq\frac14t^{2H},\qquad
 \frac{\|\mathcal R_t\|_2^2}{t^{\alpha_H}}
       \geq\frac14t^{-(\alpha_H-2H)}\longrightarrow+\infty.
\]
This is a complete quantitative exclusion of the proposed squared
remainder order. Given any real $C$ and $\delta>0$, if $C\leq0$
any $0<t<\min(1,\delta)$ violates
$\|\mathcal R_t\|_2^2\leq Ct^{\alpha_H}$. If $C>0$, choose such a
$t$ additionally satisfying $Ct^{\alpha_H-2H}<1/4$; the same
lower bound gives a strict violation. No limiting equality or
upper asymptotic for the full fractional forecast is needed.
\par\smallskip\begingroup\small\raggedright
\noindent\textit{Source correspondence:}
\code{FractionalUnifiedRemainderDiagnostic.initial_sourceSpan},
\code{initial_sourceProjection}, \code{initialLeading_zero},
\code{claimed_exponent_gt}, \code{no_claimed_remainder_bound}.
\par\endgroup

At $H=1/2$, the forecast is simply $F_t=B(t)$ for $0\leq t\leq1$,
so the actual squared remainder is $t$, whereas $\alpha_H=2$.
This Brownian special case identifies the distinction in the
leading term: $Q_0u_1=0$ is a projection of a deterministic source
direction onto the initially observed source span. It is not the
conditional expectation of the random gradient field. For the
linear terminal target that gradient is the deterministic vector
$u_1$, whose conditional expectation is still $u_1$.
The failure concerns precisely the displayed projected-source
scalar leading term and its claimed order. It does not refute
correctly formulated Brownian or fractional stochastic expansions,
which require their own integrands, domains and remainder norms.

\section{Inverse-number smoothing preserves the obstruction}
Define on the one-coordinate terminal space the inverse-number potential
\begin{equation}\label{eq:potential}
 \U f=\sum_{n\geq1}\frac{a_n}{n}W_n(u),\qquad
 f=\sum_{n\geq0}a_nW_n(u).
\end{equation}
It removes the mean. In the full Gaussian library this is the bounded
inverse-number operator constructed as the divergence recovery map
followed by its Hilbert adjoint; its action on each positive chaos is
$1/n$. The diagonal formula \eqref{eq:potential} is its exact
restriction, not a different regularizer.

\begin{proposition}\label{prop:potential}
$\U$ is bounded, self-adjoint, nonnegative and compact on
$L^2(\sigma(X(u)),P)$. Each positive coefficient is nonzero after
smoothing exactly when it was nonzero before smoothing; in particular
the least positive Hermite rank is preserved. Moreover
\begin{equation}\label{eq:potentialenergy}
 \|\E[\U f\mid\mathcal A]\|_2^2
 =\sum_{n\geq1}\frac{a_n^2}{n^2}n!v^n.
\end{equation}
For fractional terminal $f$ and any $m\geq1$ with $a_m\ne0$,
\begin{equation}\label{eq:potentiallower}
 \|P_t\U f\|_2^2\geq\frac{a_m^2m!}{m^24^m}t^{\kappa_Hm}.
\end{equation}
Every higher-power increment bound as in \eqref{eq:nohigher} therefore
fails for $\U f$ as well. Centering of the original $f$ is unnecessary.
\end{proposition}
\begin{proof}
In the orthonormal basis $W_n(u)/\sqrt{n!}$, $\U$ is diagonal with
eigenvalues $0,1,1/2,1/3,\ldots$. These real nonnegative multipliers
prove boundedness, self-adjointness and nonnegativity. Truncating after
rank $N$ gives a finite-rank operator with norm error $1/(N+1)$;
hence $\U$ is compact. Its positive coefficients are $a_n/n$, so
positive ranks are preserved. Apply Theorem~\ref{thm:energy} to this
new terminal function to obtain \eqref{eq:potentialenergy}, and
Proposition~\ref{prop:variance} to get \eqref{eq:potentiallower}.
The latter has positive constant and $P_0\U f=0$; the same elementary
power comparison proves the increment obstruction.
\end{proof}

This compactness is on a \emph{single} Gaussian observation space. It
does not imply compactness of the full Gaussian potential on an
infinite-dimensional source: on first chaos the potential acts as the
identity, and an infinite orthonormal family has no norm-convergent
subsequence. Nor does compact smoothing in Hermite rank imply better
physical-time regularity of conditional forecasts. The obstruction
retains its positive-rank power because the corresponding coefficient
is divided by a fixed nonzero integer.

\section{A concrete full-past versus current-value witness}
The Hilbert projection in the prediction formula must use the correct
observation family. An explicit example already distinguishes the full
past from the present coordinate.

\begin{proposition}\label{prop:memory}
Let $H=1/4$ and define
\[
 F=\E[B_4^H\mid\mathcal F_2^H],\qquad
 C=\E[B_4^H\mid\sigma(B_2^H)].
\]
Then $F\ne C$ in $L^2$, and $\|F\|_2^2>\|C\|_2^2$.
\end{proposition}
\begin{proof}
The exact covariance formula gives
\[
 R_H(4,2)=1,\quad R_H(4,1)=\tfrac12(3-\sqrt3),\quad
 R_H(2,1)=\tfrac12\sqrt2,\quad R_H(2,2)=\sqrt2.
\]
The one-observation projection is $C=B_2^H/\sqrt2$, so
$\langle C,B_1^H\rangle=1/2$. Since $B_1^H$ is in the full past,
$\langle F,B_1^H\rangle=R_H(4,1)=(3-\sqrt3)/2$, a different number.
Thus $F\ne C$. The one-observation subspace is contained in the
full-past $L^2$ subspace. Both are orthogonal projections of $B_4^H$;
their projection identities give $F-C\perp C$. Pythagoras yields
$\|F\|_2^2=\|C\|_2^2+\|F-C\|_2^2>\|C\|_2^2$.
\end{proof}

This witness does not assert that fractional Brownian motion is a new
non-Markovian example. It records exactly why replacing a full-past
projection with the current-coordinate projection changes the theorem.

\section{Raw quartic targets and cross-covariance of residuals}
The rank-energy identity concerns one transformed coordinate. With several
targets, their prediction errors need not be mutually orthogonal. We now
give exact formulas and an actual Gaussian example showing why diagonal
error estimates cannot be added without retaining the cross terms.

For an arbitrary direction $a\in\Hh$, not necessarily unit, write
$s_a=\|a\|^2$ and $p_a=P_Ma$. Define the centered raw fourth power
\[
 F_a=X(a)^4-3s_a^2=W_4(a)+6s_aW_2(a),
 \quad P_a=\E[F_a\mid\mathcal A],\quad R_a=F_a-P_a.
\]
The Wick expansion follows by the explicit identities
$W_2(a)=X(a)^2-s_a$ and
$W_4(a)=X(a)^4-6s_aX(a)^2+3s_a^2$.
In particular a raw fourth power has both second and fourth chaos.
It is not a pure rank-four target.

\begin{proposition}[Quartic prediction and error]\label{prop:quartic}
The actual conditional prediction is
\begin{align*}
 P_a&=W_4(p_a)+6s_aW_2(p_a)\\
 &=X(p_a)^4+6(s_a-\|p_a\|^2)X(p_a)^2
                   +3\|p_a\|^4-6s_a\|p_a\|^2.
\end{align*}
Its full and retained energies are
\[
 \|F_a\|_2^2=96s_a^4,\qquad
 \|P_a\|_2^2=24\|p_a\|^8+72s_a^2\|p_a\|^4.
\]
If $s_a>0$ and $q=\|p_a\|^2/s_a$, then
\begin{equation}\label{eq:quarticfraction}
 \frac{\|R_a\|_2^2}{\|F_a\|_2^2}
 =1-\frac34q^2-\frac14q^4.
\end{equation}
Equivalently, for the missing linear fraction $r=1-q$, this is
$(240r-216r^2+96r^3-24r^4)/96$.
\end{proposition}
\begin{proof}
The conditional Wick identity from Lemma~\ref{lem:conditional} is valid for
any direction by the same orthogonal split, or by scaling. It gives the
first prediction formula, and expansion of $W_4,W_2$ gives the second.
Lemma~\ref{lem:hermite} makes second and fourth chaoses orthogonal and gives
$\|W_n(a)\|_2^2=n!s_a^n$. Thus
\[
 \|F_a\|_2^2=24s_a^4+36s_a^2(2s_a^2)=96s_a^4.
\]
The same computation at $p_a$, retaining the coefficient $6s_a$, proves
the prediction energy. Since conditional expectation is orthogonal
projection, $R_a\perp P_a$ and
$\|R_a\|_2^2=\|F_a\|_2^2-\|P_a\|_2^2$. Dividing gives
\eqref{eq:quarticfraction}; substituting $q=1-r$ and expanding proves
the final polynomial. The second-chaos contribution is exactly the term
that a pure fourth-chaos factor would miss.
\end{proof}
Source module: \code{GaussianQuarticConditionalVariance}.
\par\noindent Declarations: \code{prediction_wick},
\code{prediction_coeFn}, \code{prediction_norm_sq},
\code{variance_fraction} and \code{residual_energy}.

\begin{proposition}[Residual covariance]\label{prop:rescov}
For any $a,b\in\Hh$, put $c=\langle a,b\rangle$ and
$d=\langle p_a,p_b\rangle$. Then
\begin{align*}
 \langle F_a,F_b\rangle&=72s_as_bc^2+24c^4,\\
 \langle P_a,P_b\rangle&=72s_as_bd^2+24d^4,\\
 \langle R_a,R_b\rangle
 &=72s_as_b(c^2-d^2)+24(c^4-d^4).
\end{align*}
The last expression need not be nonnegative, even when $c=0$.
\end{proposition}
\begin{proof}
Expand the second/fourth Wick expressions and apply
Lemma~\ref{lem:hermite}. Cross-chaos terms vanish; the remaining terms
have coefficients $36s_as_b\,2!$ and $4!$. This proves the first two
identities. Each residual is orthogonal to every $\mathcal A$-measurable
$L^2$ variable, including both predictions. Consequently
$\langle F_a,F_b\rangle=\langle P_a,P_b\rangle+\langle R_a,R_b\rangle$,
which proves the third identity without an assumption about the signs of
cross-covariances.
\end{proof}
Source module: \code{GaussianQuarticResidualCovariance}.
\par\noindent Declarations: \code{quartic_covariance},
\code{prediction_covariance}, \code{residual_prediction_orthogonal}
and \code{residual_covariance}.

\begin{theorem}[Independent targets with negative residual covariance]
Let $G_0,G_1$ be independent standard Gaussian random variables, observe
only $G_0$, and put
\[
 F_+=(G_0+G_1)^4-12,\qquad F_-=(G_0-G_1)^4-12.
\]
The targets $F_+,F_-$ are centered and independent. Their predictions are
the same variable
\[
 P=\E[F_+\mid G_0]=\E[F_-\mid G_0]=G_0^4+6G_0^2-9.
\]
For $R_\pm=F_\pm-P$, one has
\[
 \langle R_+,R_-\rangle=-312,\quad
 \|R_+\|_2^2=\|R_-\|_2^2=1224,\quad
 \|R_++R_-\|_2^2=1824.
\]
The unexplained fraction for the sum of the targets is $19/32$, strictly
below the individual fraction $51/64$.
\end{theorem}
\begin{proof}
Use orthonormal directions $e_0,e_1$ with $X(e_i)=G_i$, take
$a=e_0+e_1$, $b=e_0-e_1$, and $M=\operatorname{span}\{e_0\}$.
Then $s_a=s_b=2$, $c=0$, and $p_a=p_b=e_0$, so $d=1$.
The two linear coordinates are jointly Gaussian with covariance zero:
their joint characteristic function is
$\exp(-(2t^2+2s^2)/2)$, a product. Thus they are independent, as are their
fourth-power functions. Each fourth moment is $3\cdot2^2=12$,
so centering is correct.

Proposition~\ref{prop:quartic} gives the printed common predictor and
\[
 \|F_\pm\|_2^2=96\cdot2^4=1536,\qquad
 \|P\|_2^2=24+72\cdot2^2=312.
\]
Hence each residual has energy $1536-312=1224$.
Proposition~\ref{prop:rescov} gives
$\langle R_+,R_-\rangle=0-312=-312$. Expanding the squared norm of the
residual sum yields $1224+1224-624=1824$.
The full sum has variance $1536+1536=3072$ by independence, while its
predictor $2P$ has variance $4\cdot312=1248$.
Its unexplained fraction is therefore
$1-1248/3072=19/32$, compared with $1-312/1536=51/64$ for each target.
Since $19/32=38/64<51/64$, the inequality is strict.
\end{proof}
The mapped source additionally identifies the actual residual representatives
\[
 R_\pm=(G_0\pm G_1)^4-G_0^4-6G_0^2-3
\]
and their integral product $-312$. This agrees with the projection argument
above; it is not a covariance assigned to formal abstract energy variables.
See \code{GaussianQuarticResidualCovariance.residual_plus_coeFn},
\code{residual_minus_coeFn}, \code{example_residual_integral},
\code{example_sum_fraction} and \code{sum_fraction_below_individual}.

The inference to avoid is an additive lower bound on a block's unexplained
energy from its diagonal unexplained energies alone. The residuals are
orthogonal to the observation space, but they are not thereby orthogonal
to each other. Nor is their conditional covariance asserted negative at
every observation value: the theorem concerns their actual integrated
$L^2$ covariance. This is a concrete classical projection diagnostic, with
no new law of Gaussian prediction or multi-target error bound claimed.



\subsection*{Exact finite-block retained variance}
The cross terms can be retained exactly for any finite block of target
directions $(u_i)_{i\in J}$. Define covariance matrices and diagonal data
\[
 C_{ij}=\langle u_i,u_j\rangle,\quad
 B_{ij}=\langle P_Mu_i,P_Mu_j\rangle,\quad
 c_i=C_{ii},\quad b_i=B_{ii},\quad r_i=c_i-b_i\geq0.
\]
Orthogonal projection gives $C=B+A$, where
$A_{ij}=\langle u_i-P_Mu_i,u_j-P_Mu_j\rangle$ and $A_{ii}=r_i$.
Thus $A,B,C$ are positive semidefinite covariance matrices, but their
individual off-diagonal entries need not be positive.
For a real matrix $D$ and real diagonal lists $d,e$, use
\[
 T_1(D)=\sum_{i,j}D_{ii}D_{jj}D_{ij}^2,\qquad
 T_2(D)=\sum_{i,j}D_{ij}^4,\qquad
 \mathcal W(d,e;D)=\sum_{i,j}d_ie_jD_{ij}^2.
\]
Consider the actual Gaussian block, with its normalization fixed,
\[
 \mathcal V=\frac1{24}\sum_{i\in J}F_{u_i},
 \qquad \widehat{\mathcal V}=\E[\mathcal V\mid\mathcal A].
\]
Both are centered square-integrable variables, since the block is finite
and each Gaussian target has finite moments of every order.
Proposition~\ref{prop:rescov}, summed over all ordered pairs, gives
\begin{align}\label{eq:blockquartic}
 \|\mathcal V\|_2^2
 &=\frac{72T_1(C)+24T_2(C)}{576},\\
 \|\widehat{\mathcal V}\|_2^2
 &=\frac{72\mathcal W(c,c;B)+24T_2(B)}{576}.\nonumber
\end{align}
The denominator is $24^2$; no independent-site approximation is used.
The retained second-chaos coefficient still contains the \emph{total}
variance list $c$, not just the retained list $b$.

Expanding $c=b+r$ and using symmetry of $B$ yields the full three-term
formula
\begin{align}\label{eq:blockquarticsplit}
 \|\widehat{\mathcal V}\|_2^2
 ={}&\frac{72T_1(B)+24T_2(B)}{576}\\
 &+\frac18\mathcal W(r,r;B)
   +\frac14\mathcal W(b,r;B).\nonumber
\end{align}
Indeed bilinearity gives
$\mathcal W(b+r,b+r;B)=\mathcal W(b,b;B)
+\mathcal W(r,r;B)+2\mathcal W(b,r;B)$;
$\mathcal W(b,b;B)=T_1(B)$, and $72/576=1/8$.
All three terms in \eqref{eq:blockquarticsplit} are nonnegative,
because $b_i,r_i\geq0$ and the matrix entries are squared or raised
to the fourth power. The first term alone is the variance of
$24^{-1}\sum_i\big(X(P_Mu_i)^4-3b_i^2\big)$.
It generally undercounts the actual retained block variance: the other
two terms are the precise missing corrections.

Equivalently, writing $Y_i=X(P_Mu_i)$, the actual prediction is
\[
 \widehat{\mathcal V}
 =\frac1{24}\sum_i
 \big\{Y_i^4-3b_i^2+6r_i(Y_i^2-b_i)\big\}.
\]
This follows by the already proved quartic conditional formula. A term
involving interior variance can survive conditioning through a retained
quadratic observable. Orthogonality of a single residual to the retained
space does not remove residual cross covariances between different
sites. In particular the earlier explicit two-target example still
prevents replacing a block residual norm by its diagonal sum.
All these statements concern an actual Gaussian reference. They do not
supply an interacting Gibbs entropy inequality, entropy tensorization,
or a rate uniform in block size.
\par\smallskip\begingroup\small\raggedright
\textit{Source correspondence:} \code{GaussianQuarticBlockVariance};
\code{block_norm_sq}, \code{block_prediction_condExp},
\code{block_prediction_norm_sq}, \code{block_prediction_split}.
\par\endgroup

\subsection*{Shifted scalar fourth powers: a genuine variance minimum}
For a deterministic shift $b\in\R$ and a centered Gaussian $Z$ of variance
$r\geq0$, set
\[
 U_b=(b+Z)^4-(b^4+6b^2r+3r^2).
\]
The subtracted expression is its actual mean. If $r=0$, $Z=0$ almost
surely and $U_b=0$ for every shift. For $r>0$ use the monic variance-$r$
Hermite polynomials $H_j(Z;r)$, characterized by
$\exp(tZ-rt^2/2)=\sum_j H_j(Z;r)t^j/j!$.
Direct polynomial expansion gives
\[
 U_b=H_4+4bH_3+6(b^2+r)H_2+4b(b^2+3r)H_1.
\]
Here $H_1=Z$, $H_2=Z^2-r$, $H_3=Z^3-3rZ$ and
$H_4=Z^4-6rZ^2+3r^2$; substitution verifies both centering and the
identity. Their orthogonality follows from the same Gaussian generating
function calculation as Lemma~\ref{lem:hermite}, with squared norms
$j!r^j$. Therefore
\begin{align}\label{eq:shiftquartic}
 \E U_b^2
 &=24r^4+16b^2(6r^3)+36(b^2+r)^2(2r^2)
       +16b^2(b^2+3r)^2r\\
 &=16b^6r+168b^4r^2+384b^2r^3+96r^4.\nonumber
\end{align}
Every shift-dependent term is nonnegative, and for $r>0$, $b\ne0$,
the term $384b^2r^3$ is strictly positive. Thus the variance has its
unique minimum $96r^4$ at zero shift when the residual variance is
positive. This is a valid positive local statement, rather than an
assumption about how all conditional errors combine.

For the actual Gaussian orthogonal split $X(u)=X(P_Mu)+Z$, the residual
$Z=X(u-P_Mu)$ is independent of the entire observed sigma-algebra and
has variance $r=\|u-P_Mu\|^2$. Consequently the conditional variance
of $X(u)^4$ is the polynomial \eqref{eq:shiftquartic} with
$b=X(P_Mu)$. This identifies an actual conditional variance: integrate
only the independent Gaussian residual, first for bounded product tests
of the observed variables and then by Gaussian polynomial integrability.
For one observed coordinate, it is also the literal scalar Gaussian
fiber formula. A minimum at a deterministic shift proves neither a
positive lower bound for every block residual nor decay of boundary
information under resampling.
\par\smallskip\begingroup\small\raggedright
\textit{Source correspondence:} \code{GaussianShiftedQuarticVariance};
\code{shiftedCentered_coeFn}, \code{shiftedCentered_norm_sq},
\code{variance_minimum}, \code{variance_strict},
\code{fiber_raw_mean}, \code{fiber_variance}.
The full-observation conditional interpretation is a labelled human
application of the existing independent orthogonal split; the mapped
scalar fiber statements retain their literal one-coordinate law.
\par\endgroup

\section{Past-plus-terminal observations and bracket errors}
The same Gaussian projection mechanism controls a distinct question with
anticipative coefficients. Let $W$ be standard Brownian motion under its
original law and fix $T>0$. Write
\[
 \mathcal G_s=\sigma(W_r:0\le r\le s)\vee\sigma(W_T),\qquad 0\le s<T.
\]
This includes the entire raw past together with a single terminal value.
It is not a claim of predictability for the original Brownian filtration.
For $s\le t\le T$, put $h=t-s$, $L=T-s$ and
$Q_{s,t}=(W_t-W_s)^2-h$.

\begin{proposition}\label{prop:bridgebracket}
The actual conditional projection is
\begin{equation}\label{eq:bridgebracket}
 \E[Q_{s,t}\mid\mathcal G_s]
 =\frac{h^2}{L^2}\{(W_T-W_s)^2-L\}.
\end{equation}
Consequently, for $A\in L^2(\mathcal G_s)$,
\begin{equation}\label{eq:bridgepair}
 |\E[AQ_{s,t}]|\le \frac{2h^2}{L}\|A\|_2.
\end{equation}
\end{proposition}
\begin{proof}
Set $S=W_T-W_s$ and
$\xi=W_t-W_s-(h/L)S$. These are jointly Gaussian, with
$\E S^2=L$, $\E\xi^2=h-h^2/L$ and $\E\xi S=0$.
Each increment after $s$ is independent of the whole past up to $s$.
Thus $\xi$ is independent of the joint family consisting of $S$ and that
whole past: the finite joint Gaussian characteristic functions factor,
and a monotone-class argument extends independence from cylinders to
the generated sigma algebra. This algebra is precisely $\mathcal G_s$,
since $W_s$ is already observed. Expanding
$((h/L)S+\xi)^2-h$ and integrating the independent centered $\xi$ gives
\eqref{eq:bridgebracket}. When $h=0$ the same identity is zero.

The centered Gaussian square has second moment $2L^2$, since
$\E S^4=3L^2$. Hence the projected variable has norm
$\sqrt2 h^2/L\le2h^2/L$. The product with $A$ is integrable by
Cauchy--Schwarz. Orthogonal projection transfers its pairing to the
conditional variable, and another Cauchy--Schwarz inequality gives
\eqref{eq:bridgepair}. This is the actual finite/infinite-observation
Wick projection of the mapped source, not only conditioning on the
current endpoint in place of the entire past.
\end{proof}

Use a uniform grid $t_i=iT/n$, $0\le i\le n$, and write
$h=T/n$ and $Q_i=Q_{t_i,t_{i+1}}$. The source indexes $n=N+1$;
we use $n\ge1$ throughout. Let $H_n=\sum_{k=1}^n1/k$.
\begin{theorem}\label{thm:bracketmean}
For any triangular array $A_{n,i}\in L^2(\mathcal G_{t_i})$ with
$\|A_{n,i}\|_2\le M$,
\begin{equation}\label{eq:bracketmean}
 \left|\E\sum_{i=0}^{n-1} A_{n,i}Q_i\right|
 \le\frac{2MT}{n}H_n\le\frac{4MT}{\sqrt n}\longrightarrow0.
\end{equation}
The harmonic-sum bound and its consequence
$2MT(1+\log n)/n$ are direct human refinements of the formally recorded
last bound. No uniform boundedness or differentiability of the
coefficients is required in this mean statement.
\end{theorem}
\begin{proof}
Apply \eqref{eq:bridgepair} cell by cell. Since
$T-t_i=h(n-i)$, each pairing is at most
$2Mh/(n-i)$. Summing proves the first inequality. Also
$1/k\le1/\sqrt k$ and, by integration of $x^{-1/2}$,
$\sum_{k=1}^n1/\sqrt k\le1+\int_1^n x^{-1/2}dx
=2\sqrt n-1\le2\sqrt n$.
This gives the formally recorded estimate. Integration of $1/x$ gives
$H_n\le1+\log n$, proving the stated refinement and convergence.
Each expectation is legitimate by the pairing integrability already
proved, and the sum is finite.
\end{proof}

Uniform bounds on the separate coefficient $L^2$ norms do not replace
the observation-domain hypothesis. On $T=1$, set
$A_{n,i}=nQ_i$. Then the exact Gaussian square variance gives
\[
 \|A_{n,i}\|_2^2=n^2(2/n^2)=2,\qquad
 \E\sum_{i=0}^{n-1}A_{n,i}Q_i=n^2(2/n^2)=2
\]
for every $n$. Thus the means do not tend to zero. This coefficient
uses the unresolved current increment. In particular the whole array
cannot satisfy the generated-past-plus-terminal hypothesis, since
Theorem~\ref{thm:bracketmean} would then force its constant mean two to
vanish. The source states this existential failure of the observation
domain; we do not strengthen it to failure for every individual cell.

\section{Strong bracket convergence for bounded Borel weights}
The previous result concerns a mean, not the norm of a random sum.
A separate cross-cell calculation supplies a strong statement. Let
$\varphi:\mathbb R^2\to\mathbb R$ be Borel measurable with
$|\varphi|\le1$, and define
\[
 a_i=\varphi(W_{t_i},W_T),\quad B_i=a_iQ_i,\quad
 E_n=\sum_{i=0}^{n-1}B_i.
\]
No continuity of $\varphi$ is imposed.
\begin{theorem}\label{thm:strongbracket}
The actual finite random variables satisfy
\begin{align}
 \|E_n\|_2^2
 &\le \frac{4T^2}{n}+\frac{8T^2 H_n}{n},\label{eq:strongharmonic}\\
 \|E_n\|_2^2
 &\le \frac{4T^2}{n}+\frac{16T^2}{\sqrt n}\longrightarrow0.
 \label{eq:strongformal}
\end{align}
The second estimate is the mapped formal theorem. The sharper harmonic
estimate is a human deduction retaining the same row calculation.
Thus the displayed error is $O(\sqrt{\log(n+1)/n})$ in $L^2$, without
asserting optimality, a central limit theorem or a limiting weighted
occupation integral.
\end{theorem}
\begin{proof}
Each $Q_i$ has norm $\sqrt2h\le2h$. Multiplication by a bounded
measurable variable of absolute value at most one contracts $L^2$,
so $\|B_i\|_2\le2h$. The diagonal sum is therefore at most
$4nh^2=4T^2/n$.

For $i<j$, put $C_{ij}=Q_i a_i a_j$. Both endpoints of the $i$th
increment are at most $t_j$. The variables $a_i,a_j$ are measurable
with respect to $\mathcal G_{t_j}$, and so is $Q_i$. Thus
$C_{ij}\in L^2(\mathcal G_{t_j})$ and
$\|C_{ij}\|_2\le2h$. The exact cross pairing is
\[
 \E[B_iB_j]=\E[Q_j C_{ij}].
\]
This product is integrable: $Q_j,C_{ij}\in L^2$. Equivalently the
original $B_iB_j$ is integrable by the same inequality. No independence
of the weighted cells is assumed.

Define the later-time row $C_j=\sum_{i<j}C_{ij}$. It belongs to
$L^2(\mathcal G_{t_j})$ and satisfies
$\|C_j\|_2\le2jh\le2T$. Expanding the exact finite sum energy yields
\[
 \|E_n\|_2^2=\sum_i\|B_i\|_2^2+2\sum_j\E[Q_jC_j].
\]
Apply the first bound in \eqref{eq:bracketmean} to the row array with
$M=2T$. The absolute value of its sum of pairings is at most
$4T^2H_n/n$. Combining with the diagonal proves
\eqref{eq:strongharmonic}. Using $H_n\le2\sqrt n$ gives
\eqref{eq:strongformal}. Its right side tends to zero; taking square
roots gives the strong convergence. Finally $H_n\le1+\log n$ supplies
the printed rate refinement. All bounded products are taken in the
actual Gaussian probability space, rather than in an invented
coefficient pairing.
\end{proof}

For any finite bound $|\varphi|\le K$, multiplication and rescaling
by $K>0$ multiply these squared bounds by $K^2$; $K=0$ is immediate.
The formal source uses one fixed function $\varphi$ and the normalized
bound one. This rescaling and the harmonic refinement are explicit
classical deductions. The theorem addresses the centered bracket error
alone; it does not assert that
$\sum_i h\varphi(W_{t_i},W_T)$ converges to a time integral for an
arbitrary Borel function.

Anticipative variation theory has extensive classical antecedents.
Yoshida's primary introduction \cite{Yoshida} treats variation sums
with random weights, explicitly allows nonpredictable weights, and
sets up a substantially stronger fluctuation/expansion problem with
Malliavin--Sobolev regularity. We use this as context, not as a theorem
covering every bounded Borel weight here. The complete elementary
Gaussian projection and later-row proof above has a narrower target,
with explicit hypotheses and finite errors. No new anticipative
stochastic integral or first variation convergence principle is claimed.

\section{Recovery requires both sample and source projection}
We now distinguish the Hermite rank of a target from the dimension of
the observation space $M$. Work in the full Gaussian space
$\mathcal G=L^2(\sigma(X(h):h\in\Hh),P)$, with no additional random
factor, and in its vector counterpart $\mathcal V=L^2(P;\Hh)$ over
the same sigma-algebra. Let $Q=\E[\,\cdot\mid\mathcal A]$ on scalars,
let $\widehat Q$ be the Bochner conditional expectation on vectors,
and let $\widehat P$ act pointwise by the source projection $P_M$.
On a tensor $af$, with $a\in\Hh$, $f\in\mathcal G$,
\[
 \widehat Q(af)=aQf,\qquad \widehat P(af)=(P_Ma)f.
\]
Finite sums of such tensors are dense in $\mathcal V$. The two
operators are commuting orthogonal projections, first on these tensors
and then by continuity. Their product $J=\widehat P\widehat Q$ is
therefore the orthogonal projection onto the fields that are both
$\mathcal A$-measurable and $M$-valued.

Let $D$ denote the original closed Gaussian gradient, obtained by
closing the cylindrical polynomial gradient, and let $\delta=D^*$ be
its maximal Hilbert adjoint. Write $P_q$ for orthogonal projection to
chaos $q$. With the positive number-operator convention $N=\delta D$,
define $N^{-1}$ to be zero on constants and multiplication by $1/q$
on chaos $q\geq1$. Its associated recovery is
\begin{equation}\label{eq:Rdefinition}
 Rf=\sum_{q\geq1}\frac1q DP_qf,\qquad
 \|Rf\|_{\mathcal V}^2=\sum_{q\geq1}\frac{\|P_qf\|_2^2}{q}
 \leq\|f\|_2^2.
\end{equation}
This bounded map is defined for every Gaussian $L^2$ target; it does
not require $f\in D(D)$.

We give the graph justification rather than treating the notation
$DN^{-1}$ as a domain assertion. Finite tensor Hermite polynomials
obey $\delta Df_q=qf_q$ and
$\langle Df_q,Dg_r\rangle=\mathbf1_{q=r}q\langle f_q,g_q\rangle$.
These follow by one-coordinate Hermite differentiation, its Gaussian
integration-by-parts adjoint, and orthogonality of the product basis.
Closing these finite identities first in a fixed chaos gives
$P_qf\in D(D)$, the same gradient norm and the original adjoint
graph pair $(DP_qf,qP_qf)$. Equation \eqref{eq:Rdefinition} follows by
orthogonal summation of the gradient fields. The finite recovery sums
and their divergences converge respectively to $Rf$ and $f-\E f$;
closedness of $\delta$ proves
\begin{equation}\label{eq:Rgraph}
 Rf\in D(\delta),\qquad \delta Rf=f-\E f.
\end{equation}
Every such recovery field is orthogonal to $\ker\delta$: for a
divergence-free $v$, adjointness gives
$\langle v,DP_qf\rangle=\langle\delta v,P_qf\rangle=0$.
Pass through the convergent series. Any other field with the same
divergence differs from $Rf$ by a kernel element, so Pythagoras proves
that $Rf$ is the unique minimum-norm preimage.
\par\smallskip\begingroup\small\raggedright
\textit{Source correspondence:} \code{GaussianDivergenceRecovery.recovery_graph},
\code{recovery_inverse_energy_hasSum}, \code{recovery_orthogonal_kernel}.
\par\endgroup

\begin{theorem}[Joint observation/recovery identity]\label{thm:jointR}
For every $f\in\mathcal G$,
\[
 JRf=RQf.
\]
In particular $RQf$ is both observed and $M$-valued, and is the
minimum-norm original divergence preimage of $Qf-\E f$. Projection
of the sample variable alone generally fails: if $p=P_Mu$, then
for $q\geq1$,
\begin{align}
 \widehat Q R W_q(u)-RQW_q(u)&=(u-p)W_{q-1}(p),\label{eq:Rdefect}\\
 \|\widehat Q R W_q(u)-RQW_q(u)\|_{\mathcal V}^2
  &=(q-1)!\|u-p\|^2\|p\|^{2(q-1)}.\label{eq:Rdefectenergy}
\end{align}
At $q=1$ the defect is the deterministic field $u-p$.
\end{theorem}
\begin{proof}
The Wick derivative is $DW_q(u)=q\,uW_{q-1}(u)$, so
$RW_q(u)=uW_{q-1}(u)$. The conditional identity already proved in
Section 2 gives $QW_q(u)=W_q(p)$. Consequently
\[
 JRW_q(u)=pW_{q-1}(p)=RQW_q(u),\qquad
 \widehat Q RW_q(u)=uW_{q-1}(p).
\]
Subtracting proves \eqref{eq:Rdefect}; its norm is the product of
the deterministic squared norm and the Wick energy
$(q-1)!\|p\|^{2(q-1)}$, proving \eqref{eq:Rdefectenergy}.
The constant case has zero recovery on both sides.

Pure Wick powers span a dense subspace of $\mathcal G$. To see the
needed spanning statement, expand
$W_q(\sum_j t_ju_j)$ as a polynomial in the parameters: its coefficients
give every homogeneous Wick product of $q$ Gaussian directions.
Polarization expresses each coefficient as a finite linear combination
of values of these pure powers. Tensor Hermite polynomials are products
of Wick powers in orthogonal directions and have exactly these
homogeneous Wick products. Finite-dimensional product Hermite bases
are complete by the one-coordinate completeness proof in Section 2
and density of finite sums of product functions. Finally choose a
countable dense orthonormal source family. Gaussian $L^2$ limits of
its linear coordinates generate all $X(h)$ modulo null sets, and
finite-coordinate cylindrical functions are dense in their generated
$L^2$ space by the monotone-class approximation of indicator functions.
This proves the asserted density.

Both $JR$ and $RQ$ are bounded, so their equality on pure Wick powers
extends to all $\mathcal G$. The range property follows from the
definition of $J$; \eqref{eq:Rgraph} applied to $Qf$ and $\E Qf=\E f$
gives its divergence. The minimum-norm conclusion was proved above.
\end{proof}
\par\smallskip\begingroup\small\raggedright
\textit{Source correspondence:} \code{GaussianWickRecoveryPrediction.prediction_recovery_defect},
\code{prediction_recovery_defect_energy};
\code{GaussianJointSourceProjection.jointConditional_fixed_iff};
\code{GaussianJointRecoveryIntertwining.joint_recovery_operator}.
\par\endgroup
These are source/sample constraints at a fixed observation family.
They do not by themselves identify an adapted physical-time stochastic
integrand or a Clark--Ocone representation for a different observed process.

\section{A stationary factor need not decrease defect energy}
The preceding source/sample identities use specific compatible Gaussian
operators. General stationary observation factors require their own
energy argument. Here is the classical positive mechanism. On a real
Hilbert space let $T$ be a bounded symmetric contraction and let $Q$ be
the orthogonal projection onto a closed $T$-invariant subspace $V$.
Symmetry makes $V^\perp$ invariant as well: if $r\perp V$ and $v\in V$,
$\langle Tr,v\rangle=\langle r,Tv\rangle=0$. Thus $TQ=QT$.
Writing $f=Qf+r$ gives the exact splitting
\begin{equation}\label{eq:defectsplit}
 \langle(I-T)f,f\rangle
 =\langle(I-T)Qf,Qf\rangle+\langle(I-T)r,r\rangle.
\end{equation}
Since $\langle Tr,r\rangle\leq\|Tr\|\|r\|\leq\|r\|^2$,
the residual term is nonnegative. Hence $Q$ decreases defect energy.
The isometric-factor version in the frozen library uses $Q=JJ^*$,
$TJ=JS$ and the same orthogonal decomposition; it is not a separate
probabilistic assumption. In particular, for each $t>0$ the conclusion
applies to $t^{-1}(I-T_t)$ when $T_t$ is a symmetric contraction
preserving the factor. Uniform upper bounds on these defect energies
then transfer to the factor, as do existing finite limiting energies.
No unbounded-generator domain is inferred just from one bounded defect.

Stationarity and factor invariance alone do not suffice. Take the
uniform probability law on $\{0,1,2\}$ and the Markov matrix
\[
 T=\begin{pmatrix}0&1&0\\1/2&0&1/2\\1/2&0&1/2\end{pmatrix},
 \qquad Qf=\left(f_0,\frac{f_1+f_2}{2},\frac{f_1+f_2}{2}\right).
\]
Every row and column of $T$ sums to one, all entries are nonnegative,
and $Q$ is exactly conditional expectation on the partition
$\{\{0\},\{1,2\}\}$. Jensen's inequality and column sums prove
$\|Tf\|_{L^2}^2\leq\|f\|_{L^2}^2$. The range of $Q$ is $T$-invariant,
since
\[
 T(a,b,b)=(b,(a+b)/2,(a+b)/2).
\]
Its genuine coarse Markov matrix is therefore
$\left(\begin{smallmatrix}0&1\\1/2&1/2\end{smallmatrix}\right)$,
with stationary masses $1/3,2/3$.
Nevertheless the original $T$ is not symmetric:
$T_{01}=1\ne1/2=T_{10}$. For
\[
 f=(2,-1/4,-7/4),\qquad Qf=(2,-1,-1),
 \qquad \mathcal E(f)=\frac13\sum_{j=0}^2(f_j-(Tf)_j)f_j,
\]
direct evaluation gives
\[
 Tf=(-1/4,1/8,1/8),\quad TQf=(-1,1/2,1/2),\quad
 \mathcal E(f)=21/8<3=\mathcal E(Qf).
\]
Also $QTf=Tf\ne TQf$. Thus the symmetric splitting cannot be
imported merely because the retained subspace is invariant. Both
vectors are centered in the actual stationary law. Defect energy is
nonnegative here, by the same contraction bound, even though projection
can increase it. The conventional reversible Dirichlet-form setting
is recalled in Levin--Peres~\cite{LP}, Section 13.2. The example
tests the stated generalization; no new Markov-energy principle or
novelty for its finite mechanism is asserted.
\par\smallskip\begingroup\small\raggedright
\textit{Source correspondence:} \code{SymmetricFactorEnergy.quadratic_energy_split},
\code{SymmetricDefectEnergy.defect_energy_contraction};
\code{NonreversibleFactorCounterexample.retained_subspace_invariant},
\code{energy_witness}, \code{energy_projected_witness}.
\par\endgroup

\section{Finite observed dimension and compact recovery}
\begin{theorem}[Two different compactness classifications]\label{thm:compactobs}
For an arbitrary observation index set $I$,
\[
 Q\text{ is compact}\ \Longleftrightarrow\ M=\{0\},\qquad
 RQ\text{ is compact}\ \Longleftrightarrow\ \dim M<\infty.
\]
If $m=\dim M<\infty$ and $S_N=\sum_{q=0}^NP_q$, then
$RS_NQ$ has finite rank at most $(N+1)^m$, and
\begin{equation}\label{eq:Rtail}
 \|RQ-RS_NQ\|\leq(N+1)^{-1/2}.
\end{equation}
The combined rank bound is an elementary consequence of the finite
Hermite construction below; the source gives the corresponding
per-chaos bounded-Hermite-space dimension and the recovery tail estimate.
\end{theorem}
\begin{proof}
If $M\ne0$, choose a unit $u\in M$. The normalized Wick variables
$W_q(u)/\sqrt{q!}$, $q\geq0$, are orthonormal and fixed by $Q$.
They have no norm-convergent subsequence, so $Q$ is not compact.
If $M=0$, every observation is zero, $Qf=\E f$ is a rank-one map,
and $RQ=0$.

Suppose $\dim M=m<\infty$, and choose an orthonormal basis
$u_1,\ldots,u_m$. Their Gaussian coordinates are independent standard
normals. The observation $L^2$ space has product Hermite basis
\[
 H_\alpha=\prod_{j=1}^m
       \frac{\operatorname{He}_{\alpha_j}(X(u_j))}{\sqrt{\alpha_j!}},
 \qquad \alpha\in\mathbb N^m.
\]
This is the same observed space even if infinitely many labels or a
singular redundant covariance matrix describe it: source-norm convergence
and Gaussian $L^2$ limits identify the generated sigma-algebras modulo
null sets. The product $H_\alpha$ belongs to chaos $|\alpha|$.
Thus $S_NQ$ takes values in the span of the finitely many $H_\alpha$
with $|\alpha|\leq N$. Each $\alpha_j\leq N$, so there are at most
$(N+1)^m$ such indices. Applying $R$ cannot increase rank.

The chaos decomposition is preserved by $Q$, as follows on pure Wick
powers from their conditional identity and then by density. Using
\eqref{eq:Rdefinition},
\[
 \|R(I-S_N)Qf\|^2
   =\sum_{q>N}\frac{\|P_qQf\|^2}{q}
   \leq\frac1{N+1}\|Qf\|^2
   \leq\frac1{N+1}\|f\|^2.
\]
This proves \eqref{eq:Rtail}. A norm limit of finite-rank operators
is compact: their images give finite nets for the limiting image of
the unit ball, which is relatively compact in the complete target.
Hence $RQ$ is compact.

Conversely if $M$ is infinite dimensional, choose an orthonormal
sequence $u_n\in M$. The scalar inputs $X(u_n)$ have norm one,
are observed, and $RQX(u_n)=u_n$, viewed as constant fields.
The outputs are pairwise at distance $\sqrt2$, contradicting
compactness. This proves necessity and completes both classifications.
\end{proof}
\par\smallskip\begingroup\small\raggedright
\textit{Source correspondence:} \code{GaussianConditionalCompactness.conditionalProjection_compact_iff};
\code{GaussianFiniteObservationRecovery.boundedHermiteSpace_dimension},
\code{observed_chaos_block_compact}, \code{recovery_compact_iff_finite_source};
\code{GaussianRecoverySpectralCompactness.recovery_tail_operator_bound}.
\par\endgroup

Even compact one-source recovery has infinite orthonormal image energy.
For unit $u\in M$, its normalized Wick inputs satisfy
\[
 \left\|RQ\frac{W_q(u)}{\sqrt{q!}}\right\|^2=\frac1q\quad(q\geq1).
\]
Their squared image norms have divergent harmonic sum. In particular
$RQ$ is not Hilbert--Schmidt whenever $M\ne0$, including $m=1$.
Decay of each image does not give summability of all images. This is
the explicit formal harmonic-energy obstruction, not a new compactness
principle.
\par\smallskip\begingroup\small\raggedright
\textit{Source correspondence:} \code{GaussianObservedRecoveryHarmonicEnergy.observed_normalized_energy},
\code{observed_energy_not_summable}, \code{one_source_compact_divergent_energy}.
\par\endgroup

For infinite-dimensional $M$ there is a stronger exact quantitative
boundary:
\begin{equation}\label{eq:essentialone}
 \inf\{\|RQ-K\|:K:\mathcal G\to\mathcal V\text{ compact}\}=1.
\end{equation}
Indeed $X(u_n)$ above is orthonormal and hence weakly tends to zero
by Bessel's inequality. A compact operator sends this sequence to zero
in norm: otherwise a nonvanishing subsequence would have a norm
convergent image subsequence, whose limit must be zero by weak
continuity, a contradiction. Therefore
$\|(RQ-K)X(u_n)\|\geq1-\|KX(u_n)\|\to1$ and $\|RQ-K\|\geq1$.
Taking $K=0$ and $\|RQ\|\leq1$ proves \eqref{eq:essentialone}.
The frozen source states this distance for the full recovery $R$;
its observed infinite-dimensional version here is an explicitly labelled
human deduction from the same first-chaos isometry argument.
\par\smallskip\begingroup\small\raggedright
\textit{Source correspondence:} \code{GaussianRecoveryCompactApproximationGap.compact_approximation_gap},
\code{optimal_compact_error} (full recovery).
\par\endgroup

\section{An exact positive-chaos criterion for general inputs}
The preceding classification concerns a Gaussian conditional projection.
There is also a general operator statement that separates chaos-tail
decay from compactness inside each fixed chaos.
\begin{theorem}\label{thm:sectorcompact}
Let $E$ be a real normed space and $K:E\to\mathcal G$ a bounded
linear operator. Then
\[
 RK\text{ is compact}\quad\Longleftrightarrow\quad
 P_qK\text{ is compact for every integer }q\geq1.
\]
No compactness condition on the constant component is required.
\end{theorem}
\begin{proof}
Polarizing \eqref{eq:Rdefinition} gives
$R^*R=\sum_{q\geq1}q^{-1}P_q=N^{-1}$ as a bounded operator.
For fixed $q\geq1$, define the bounded extraction
$A_q=qP_qR^*:\mathcal V\to\mathcal G$. Then
$A_qR=P_q$. Thus compactness of $RK$ implies compactness of
$P_qK=A_qRK$, since bounded maps preserve relative compactness.

Conversely suppose every $P_qK$, $q\geq1$, is compact. The finite
sum $RS_NK=\sum_{q=1}^NRP_qK$ is compact. The same tail proof as
\eqref{eq:Rtail}, now applied to $Kx$, gives
\[
 \|RK-RS_NK\|\leq(N+1)^{-1/2}\|K\|\longrightarrow0.
\]
Norm closure of the compact operators proves the result. Constants
are annihilated by $R$, explaining why they do not appear in the
criterion. No finite-dimensionality of a full chaos is assumed here.
\end{proof}
\par\smallskip\begingroup\small\raggedright
\textit{Source correspondence:} \code{GaussianRecoverySpectralCompactness.sectorExtract_recovery},
\code{partial_composition_compact}, \code{recovery_composition_compact_iff}.
\par\endgroup
The theorem explains why inverse-number damping can have a small
high-chaos tail while retaining an infinite first-chaos obstruction.
It also explains finite-source compactness: each observed fixed-chaos
block has a finite Hermite range. These are classical Gaussian spectral
mechanisms. Baudoin's primary lecture \cite{BaudoinRecovery} reconstructs
the chaos decomposition and derivative-energy identity underlying
\eqref{eq:Rdefinition}; its Hermite convention includes $1/q!$ in
the polynomial, whereas the present $W_q$ is monic. The exact operator
classifications and their formal source reconstruction are presented
without a priority claim. Neither compactness nor its failure supplies
new fractional-time regularity, curvature or an observed-process
stochastic representation.



\section{Integrated history: a sufficient state need not compress recovery}
The compactness theorem concerns all observed inputs. A finite state
sufficient for one terminal target does not reduce that whole input
space to finitely many source directions. Integrated Brownian motion
gives an exact positive example, with a simultaneous failure of prediction
from its scalar current value.

Use the actual time-source space $L^2(\R,dx)$, isometrically embedded
in the Gaussian coefficient space as in the preceding constructions.
For $t\geq0$ put
\[
 b_t=1_{(0,t]},\qquad k_t(x)=(t-x)1_{(0,t]}(x),\qquad
 B_t=X(b_t),\quad J_t=X(k_t).
\]
This is the usual integrated Brownian Gaussian family. The map $r\mapsto b_r$ is continuous in $L^2$ and has norm $\sqrt r$, so it is Bochner integrable on each bounded time interval. Indeed
$\int_0^t b_r\,dr=k_t$ as a Bochner $L^2$ integral, by integration
on the finite rectangle $0\leq x,r\leq t$: at fixed $x$ its length
is $t-x$ for $0<x\leq t$, and zero otherwise. The isometry $X$
commutes with this integral, giving $J_t=\int_0^tB_r\,dr$ in $L^2$.
This elementary identification is a human reconstruction of the
formal kernel realization. For $0\leq r\leq t$, direct integration
gives
\begin{equation}\label{eq:integratedcov}
 \E[J_tJ_r]=\int_0^r(t-x)(r-x)\,dx
              =\frac{tr^2}{2}-\frac{r^3}{6},\qquad
 \E[J_r^2]=\frac{r^3}{3}.
\end{equation}

First, the integrated history recovers Brownian state in the closed
Gaussian source span. For $0\leq a<s$, let $h=s-a$. The quotient
$(k_s-k_a)/h$ equals $b_s$ on $(0,a]$, while on $(a,s]$ its
difference from $b_s$ is $-(x-a)/h$. Consequently
\begin{equation}\label{eq:integratedderivative}
 \left\|\frac{k_s-k_a}{s-a}-b_s\right\|^2
       =\frac1{h^2}\int_a^s(x-a)^2\,dx=\frac h3.
\end{equation}
The source explicitly proves the sufficient bound by $h$; the exact
$h/3$ evaluation here is a labelled classical human refinement.
Taking $a=s n/(n+1)$ gives convergence to $b_s$. The same argument
at every $r\leq s$ places $b_r$ in the closed span of
$\{k_v:0\leq v\leq s\}$. At $r=0$ the vector is zero.
Thus the Brownian state is an $L^2$-measurable function of the
integrated history up to that time, with no terminal observation added.
This is a statement about completed $L^2$ observation spaces; literal
raw sigma-algebra equality is not asserted.

Let $Q_s^J$ condition on $\sigma(J_r:0\leq r\leq s)$, and
$Q_s^{(2)}$ on $\sigma(J_s,B_s)$. For $s\leq t$, put $h=t-s$.
The exact source decomposition is
\[
 k_t=k_s+h b_s+(t-x)1_{(s,t]}(x).
\]
The last term is orthogonal to every source supported in $(0,s]$.
The forecast vector $k_s+h b_s$ lies both in the integrated-history
closed span, by \eqref{eq:integratedderivative}, and in the two-state
span. It follows that
\begin{equation}\label{eq:integratedforecast}
 Q_s^J J_t=Q_s^{(2)}J_t=J_s+hB_s,\qquad
 Q_s^J W_q(k_t)=Q_s^{(2)}W_q(k_t)=W_q(k_s+h b_s).
\end{equation}
Hermite completeness for the one-dimensional Gaussian $J_t$, and
contraction of both projections, extend the identity to
\begin{equation}\label{eq:integratedstate}
 Q_s^J f=Q_s^{(2)}f\qquad
       \bigl(f\in L^2(\sigma(J_t),P)\bigr).
\end{equation}
For $t=0$ only constants occur; otherwise rescale $k_t$ to unit norm
and apply the expansion of Section 1. Hence the entire terminal
$L^2$ space has a two-coordinate sufficient state, not just the
linear terminal value.

The scalar current value $J_s$ does not suffice. At $s=2,t=3$,
\eqref{eq:integratedcov} gives its scalar regression
$\E[J_3\mid J_2]=\tfrac74J_2$. Testing against the earlier
observation $J_1$ gives
\[
 \E[(Q_2^J J_3)J_1]=\E[J_3J_1]=\frac43,\qquad
 \E\left[\left(\frac74J_2\right)J_1\right]=\frac{35}{24}.
\]
The first equality uses orthogonality of the conditional residual to
the history. Since the two numbers differ, the conditional forecasts
differ. A further classical calculation quantifies the loss: for
$s>0,h=t-s\geq0$, the two-state prediction error has variance
$h^3/3$, the squared norm of the future kernel in the decomposition.
The scalar regression coefficient is $1+3h/(2s)$, so its error
variance, by subtracting squared covariance divided by $s^3/3$,
is $sh^2/4+h^3/3$. The extra term is strictly positive for $h>0$.

Nevertheless the whole integrated-history source span is infinite
dimensional whenever $s>0$. It contains every $b_r$, $r\leq s$,
and therefore contains the normalized indicators of the disjoint
intervals $(s2^{-j-1},s2^{-j}]$, $j\geq0$. They are orthonormal.
Their first-chaos Gaussian observables are fixed by $Q_s^J$, and
recovery maps them to the corresponding deterministic unit source
fields. These images have pairwise distance $\sqrt2$, so $RQ_s^J$
is not compact. In contrast the Gram determinant of $k_s,b_s$ is
\[
 \|k_s\|^2\|b_s\|^2-\langle k_s,b_s\rangle^2
            =\frac{s^4}{12}>0.
\]
The current-state source span has dimension two; Theorem~\ref{thm:compactobs} therefore
makes $RQ_s^{(2)}$ compact. These recovery conclusions are explicit
human applications of the already proved source-dimension criterion.
They are compatible with \eqref{eq:integratedstate}, which agrees
only on the specified terminal subspace. Neither all observed inputs
nor the full recovery operators have been identified.

Integrated Gaussian prediction is classical. Lachal's primary author
record \cite{Lachal} describes iterated Brownian integrals and their
prediction problems; only that record's stated context, not an
uninspected bridge theorem, is used here. No new state-augmentation
or Markov principle is claimed. This integrated Brownian family is
also distinct from iterated Brownian motion with a random clock.

\par\smallskip\begingroup\small\raggedright
\textit{Source correspondence:}
\code{IntegratedBrownianStatePrediction} and
\code{IntegratedBrownianHistoryReconstruction} give actual source
projections, all-terminal-$L^2$ state sufficiency and the scalar
counterexample. \code{BrownianHistoryRecoveryCompactness} constructs
the disjoint Brownian source directions; its adaptation to integrated
history uses the displayed recovered-state inclusion.
\par\endgroup

\section{A continuous observation with instantaneous information arrival}
The geometry of an observation space must be checked independently
of continuity of the observed paths. The following classical example
also explains why a raw natural filtration cannot silently be replaced
by its right-continuous enlargement. It concerns a separate Gaussian
observation, not an assertion about the IBM natural filtration.

Let $G=X(u)$ for a unit source $u\in\Hh$, so $G$ is standard normal.
For real $t$ put
\[
 Y_t=\max(t,0)G,\qquad
 \mathcal F_t=\sigma(Y_s:s\leq t),\qquad
 M_t=\mathbf1_{\{t>0\}}G.
\]
Every path of $Y$ is continuous. Nevertheless
\[
 \mathcal F_t=\begin{cases}\{\varnothing,\Omega\},&t\leq0,\\
 \sigma(G),&t>0.\end{cases}
\]
Indeed all observations before zero vanish; for positive $t$, all
past observations are functions of $G$ and $G=Y_t/t$. In particular
$\mathcal F_{0+}=\bigcap_{t>0}\mathcal F_t=\sigma(G)$ differs from
$\mathcal F_0$, since $\mathbb P(G>0)=1/2$.

\begin{proposition}\label{prop:rawarrival}
The process $M$ is a square-integrable martingale for the raw
filtration $\mathcal F$, with mean zero at every time. For every
$t>0$,
\[
 \mathbb E|M_t-M_0|^2=1.
\]
It has no modification whose paths are almost surely right-continuous
at zero. If $\Pi_t$ denotes conditional expectation onto
$\mathcal F_t$ as an operator on ambient $L^2$, then
\[
 \|\Pi_t-\Pi_0\|_{L^2\to L^2}=1\quad(t>0),
 \qquad \|Y_t\|_2=t\longrightarrow0\quad(t\downarrow0).
\]
\end{proposition}
\begin{proof}
Adaptedness and square integrability follow from the displayed
filtration and $\mathbb E G^2=1$. For $s\leq t$, if $s>0$ then
$M_s=M_t=G$ and this variable is $\mathcal F_s$-measurable.
If $s\leq0$, conditional expectation onto the trivial sigma algebra
is the scalar mean, which is zero whether $M_t=0$ or $M_t=G$.
These two cases prove the martingale property. The increment and
observation norm identities follow directly.

Let $\widetilde M$ be any modification, so
$\widetilde M_t=M_t$ almost surely for each fixed $t$.
Intersect these full-measure events only at $t=0$ and $t=1/n$,
$n\geq1$. On their common event,
$\widetilde M_0=0$ and $\widetilde M_{1/n}=G$ for every $n$.
Right continuity at zero would force $G=0$. The latter has probability
zero under the actual standard Gaussian law. Thus no right-continuous
modification exists; no uncountable intersection of exceptional sets
has been used.

For positive $t$, $\Pi_t$ is the orthogonal projection onto
$L^2(\sigma(G))$, while $\Pi_0$ projects onto constants. These
closed spaces are nested; hence $\Pi_t-\Pi_0$ is the orthogonal
projection onto the mean-zero part of $L^2(\sigma(G))$ and has
norm at most one. Its value on the unit vector $G$ is $G$, proving
the exact norm one.
\end{proof}

Completing the raw sigma algebras by null sets does not change these
conditional expectation operators: the time-zero space still contains
only almost surely constant variables. By contrast, right-continuous
enlargement inserts $G$ into the time-zero information. In that enlarged
filtration the constant process $M_t^+=G$ for $t\geq0$ is a continuous
martingale, but is not a modification of $M$ because its time-zero
value differs almost surely. Thus a regularization that changes the
initial information has changed the prediction problem.

The raw-filtration and martingale assertions are mapped to
\begin{center}\code{ContinuousObservationFiltrationObstruction}.\end{center}
The exact projection
norm and the completion/enlargement explanation are labelled classical
human deductions proved above. The mechanism itself is elementary
and established: Pitera's author lecture notes \cite{Pitera}, Section
2.3, footnote 7 on printed page 10, use $Y_t=tZ$ with a uniform
nondegenerate $Z$ to demonstrate the same filtration issue.
Gaussianity supplies our particular $L^2$ normalization, not a new
filtration principle. Neither a continuous sample path nor a vanishing
observation amplitude proves continuous arrival of information.

\subsection*{The increment energy need not be a time measure}
For this same actual raw martingale, the origin increment energy is
$\E|M_t-M_0|^2=1$ for every $t>0$. There is no countably additive
measure $\rho$ on $\R$ such that
\[
 \rho((0,t])=\E|M_t-M_0|^2\quad(t>0).
\]
Indeed $E_n=(0,1/(n+1)]$ decreases to the empty set, while
$\rho(E_n)=1$. The first interval has finite mass, so continuity
from above gives $\rho(E_n)\to\rho(\varnothing)=0$, a contradiction.
No finiteness of the whole measure is required. Placing an atom at
zero does not repair these values: zero belongs to none of $E_n$.
This conclusion concerns the specified usual half-open interval
convention and the original time-zero value.

Consequently an energy measure used to complete physical-time step
integrands is a genuine additional obligation in a raw filtration.
Square integrability and the martingale property alone do not supply
it in this generality. This does not contradict constructions under
right-continuity or other hypotheses proving the required measure.
The continuity-from-above proof is elementary; the actual source
correspondences are
\begin{center}\code{RawMartingaleEnergyMeasureObstruction}.\end{center}


\section{An accelerated Gaussian kernel limit with an initial jump}
\label{sec:instantkernel}
The preceding observation example retains information instantly. A
separate, equally classical limit example loses all state information
instantly. It concerns genuine Gaussian Markov kernels rather than
conditional projection onto a nested observation family.
Let $\gamma=N(0,1)$ and, for $u\geq0$, define
\[
 K_u(x,\cdot)=\operatorname{Law}
 \bigl(e^{-u}x+\sqrt{1-e^{-2u}}Z\bigr),\qquad Z\sim\gamma.
\]
The displayed affine map is jointly Borel in $(x,Z)$, so it defines
a measurable probability kernel. Its action preserves $C_b(\R)$
and contracts the supremum norm. For the first assertion, if
$x_j\to x$ then the affine argument converges pointwise in $Z$;
bounded convergence applies to each $f\in C_b(\R)$. The contraction
bound follows by integrating $|f|\leq\|f\|_\infty$.

Two successive transitions, with independent normals $Z_1,Z_2$, give
\[
 e^{-(u+v)}x+e^{-v}\sqrt{1-e^{-2u}}Z_1
                      +\sqrt{1-e^{-2v}}Z_2.
\]
The last two terms form a centered normal of variance
$e^{-2v}(1-e^{-2u})+1-e^{-2v}=1-e^{-2(u+v)}$.
Thus $K_{u+v}=K_uK_v$, with $K_0(x,\cdot)=\delta_x$.
Starting from an independent standard normal $x$ instead gives
variance $e^{-2u}+1-e^{-2u}=1$, proving invariance of $\gamma$.
Fubini is justified for bounded observables throughout. These are
ordinary Mehler kernels; no new kernel construction is claimed.

Set $r_n=n+1$ and $K_t^{(n)}=K_{r_nt}$ for $t\geq0$.
At fixed $n,x$ the row is weakly continuous in $t$: the affine
argument is continuous and bounded continuous tests admit dominated
convergence. For every fixed $x,t$, on the other hand,
\begin{equation}\label{eq:refreshlimit}
 K_t^{(n)}(x,\cdot)\ \Longrightarrow\ J_t(x,\cdot),\qquad
 J_t(x,\cdot)=\begin{cases}\delta_x,&t=0,\\\gamma,&t>0.\end{cases}
\end{equation}
At positive $t$ the two affine coefficients converge to zero and one,
respectively. A common $Z$ then proves the weak convergence by
bounded convergence. At zero all rows are exactly $\delta_x$.
The limit still obeys the kernel semigroup law: if either time is
zero use the identity kernel; if both are positive both compositions
are the constant kernel $\gamma$. It also preserves $\gamma$.
These algebraic properties have not supplied continuity in time.

\begin{proposition}\label{prop:refreshgenerator}
On the Banach space $C_b(\R)$ with its supremum norm, put
\[S_tf(x)=\int f(y)J_t(x,dy).\]
The norm-continuity domain at zero and
the right norm-derivative domain are both precisely the constant
functions. On that domain the derivative is zero. In particular
this right-derivative domain is closed and is not dense.
For the bounded test $f(x)=\cos x$,
\[
 \|S_tf-f\|_\infty=1+e^{-1/2}\qquad(t>0).
\]
\end{proposition}
\begin{proof}
Let $Pf=(\int f\,d\gamma)\mathbf1$. Then $P$ is a bounded
idempotent and $S_0=I$, $S_t=P$ for positive $t$.
For $t>0$ the norm difference from $f$ is the fixed number
$\|Pf-f\|_\infty$. It tends to zero exactly when $Pf=f$,
which is equivalent to $f$ being constant. Similarly the difference
quotient is $(Pf-f)/t$. A finite norm limit forces $Pf=f$ and
then that quotient vanishes. The constants are closed under uniform
limits and form a proper subspace, so are not dense. Finally the
Gaussian characteristic function gives $P\cos=e^{-1/2}$.
The bound $|e^{-1/2}-\cos x|\leq1+e^{-1/2}$ is attained at
$x=\pi$, proving the exact norm.
\end{proof}
Even pointwise temporal continuity fails at $x=0$, since
$S_0\cos(0)=1$ and $S_t\cos(0)=e^{-1/2}$ for positive $t$.
The approximating-row continuity above is weak row continuity,
not strong continuity on the full supremum-norm space $C_b$.
Classical Mehler semigroups themselves require this topological
care; see Lunardi--Pallara \cite{LunardiPallara}, Section 3.
The right norm derivative just computed is not a mixed-topology
or bounded-pointwise generator.

\subsection*{A full-past Markov realization without continuous paths}
There nevertheless is an actual coordinate probability realization.
On $\Omega=\R^{[0,\infty)}$ with its product sigma algebra take
\[
 P_x=\delta_x\ \text{at coordinate }0
           \quad\otimes\quad\bigotimes_{t>0}\gamma,
 \qquad Z_t(\omega)=\omega(t).
\]
Here the arbitrary product probability theorem is the standard
existence input. Evaluation coordinates are measurable. Finite
restrictions have the prescribed independent product laws, so
$Z_0=x$ almost surely and $Z_t\sim\gamma$ for $t>0$.
For $0\leq s<t$, $Z_t$ is independent of the entire raw past
$\mathcal H_s=\sigma(Z_r:0\leq r\leq s)$, not merely of each
individual past coordinate. To see this, independence first holds
against finite-coordinate cylinder events by the product law.
Those cylinders are a pi-system generating $\mathcal H_s$; the
class of past events satisfying the independence identity is a
lambda-system. The pi-lambda theorem proves the claim.
Consequently, for every bounded Borel $f$,
\[
 \E_x[f(Z_t)\mid\mathcal H_s]=\int f\,d\gamma
       =\int f(y)J_{t-s}(Z_s,dy)\qquad(s<t).
\]
For $s=t$ this follows from adaptedness and $J_0=I$.
Thus the full-past Markov identity holds despite the temporal jump.
There is also an explicit conditional distribution given the whole
past path. If $p\in\R^{[0,s]}$, set
$Q_{s,t}(p,\cdot)=\gamma$ for $t>s$ and
$Q_{s,s}(p,\cdot)=\delta_{p(s)}$.
This is a measurable kernel, since evaluation at $s$ is measurable.
For $t>s$ the joint law of the past path and $Z_t$ is the product
of the past law and $\gamma$, by the same cylinder argument.
For $t=s$ the future coordinate is its past evaluation.
These identities prove the conditional-distribution claim directly,
without a general disintegration theorem on an uncountable product.
The source uses real-indexed coordinates; restricting its past to
$[0,s]$ gives the version here by the same independence proof.

No realization with these one-time laws can have almost surely
right-continuous paths at zero. Indeed take $t_j=1/j$.
For a fixed start $x$, choose $\epsilon>0$ such that
$\gamma(\{y:|y-x|>\epsilon\})>0$; any finite $\epsilon>0$
works for a nondegenerate Gaussian. Then in every realization
\[
 P_x(|Z_{t_j}-x|>\epsilon)
   =\gamma(\{y:|y-x|>\epsilon\})>0
\]
for every $j$. Thus convergence to $x$ in probability fails.
Almost sure right continuity with $Z_0=x$ would imply almost sure,
hence in-probability, convergence on this countable sequence, a
contradiction. A coordinatewise modification preserves these laws,
so cannot repair the obstruction. The product construction asserts
coordinate measurability and the stated full-past laws; joint
measurability in time and sample is not assumed. A valid kernel
semigroup, a Markov conditional law and weak row limits are therefore
insufficient by themselves for a continuous-path limit or a dense
strong-generator domain. These are classical diagnostic mechanisms,
not a new general semigroup convergence theorem.

\section{Formal scope and primary comparisons}
\begingroup\raggedright
The verified baseline is
\code{bab688bb6140cd6316c80680095a914ea0d346db}.
The core correspondence is \code{GaussianWickConditionalProjection},
\code{GaussianTerminalPredictionSpectrum} and
\code{GaussianPredictionRankAsymptotic}. Actual continuous-past
identification is in \code{ContinuousFractionalPrediction}; the
two-regime obstruction is \code{FractionalDualRegimePrediction}.
The inverse-number application is \code{GaussianPotentialPredictionRank},
with the actual operator in \code{GaussianInverseNumberPotential}.
Raw quartic conditional energies and cross-target residuals are in
\code{GaussianQuarticConditionalVariance} and
\code{GaussianQuarticResidualCovariance}. The quarter-Hurst witness is
\code{FractionalCurrentObservationObstruction}.
The terminal bracket application is mapped to \code{GaussianTerminalBracketPrediction}, \code{GaussianTerminalWeightedBracket}, \code{GaussianTerminalBracketObstruction}, \code{GaussianTerminalBracketProducts} and \code{GaussianTerminalStrongBracket}. Its harmonic refinements are separately labelled human deductions. Exact declarations, original hypotheses, source hashes and human proof
reorganizations accompany the manuscript. Sections 9--12 add the
original joint source/sample recovery identity, the exact finite-source
compactness classification and the positive-chaos block criterion. The
combined finite-rank estimate and observed distance-one conclusion are
labelled human deductions. Fifty exact finite checks now pass in retained
and new groups. A further 42 exact integrated-state covariance, regression and quotient checks pass; their infinite-domain conclusions use the printed proofs. \code{verify_recovery.py} checks conditional Wick defects
through degree eight and two-coordinate derivative/recovery energies
through degree four. Infinite compactness and density arguments rest on
the printed proofs. The elementary noncompactness
of the full potential is a classical deduction from its exact first-chaos
identity, not a newly named formal theorem.

McElroy--Das's primary introduction describes the Hermite predictor and
mean-square-error mechanism for finite data; access to its complete
main-theorem text was incomplete, so no detailed theorem-number priority
comparison is claimed. Sottinen--Viitasaari's accessible primary paper,
Theorem 3.1 and Proposition 3.2, gives the conditional law and the two
small-past covariance asymptotics. Our proof retains a weaker explicit
uniform lower constant and does not assert a new fractional exponent.
The methods and central structure are classical. This is a consolidated
formalization companion with explicit domain and regularity diagnostics;
novelty remains unresolved. It neither completes the paused Bakry/fBM
campaign nor proves new curvature, temporal stochastic integration or
neural-network performance claims. The Mehler implementation retains the Michael R. Douglas adaptation
credit and Apache-2.0 notice in its original source; see
\code{LICENSES/Apache-2.0.txt} and \code{MATHEMATICAL_AUDIT.md}.
Full baseline verification evidence
was hash-checked without repeating a research build.
\par\endgroup

\clearpage
\begin{thebibliography}{9}
\bibitem{LunardiPallara} A. Lunardi and D. Pallara,
\emph{Ornstein--Uhlenbeck semigroups in infinite dimension},
Philosophical Transactions A \textbf{378} (2020), article 20190620,
Section 3. \href{https://doi.org/10.1098/rsta.2019.0620}
{doi:10.1098/rsta.2019.0620}.

\bibitem{BaudoinRecovery} F. Baudoin, \emph{Lecture 34: The Wiener chaos expansion},
primary author lecture, November 4, 2012.
\href{https://fabricebaudoin.blog/2012/11/04/lecture-34-the-wiener-chaos-expansion/}{Author lecture}.
\bibitem{MD} T. McElroy and S. Das,
\emph{Nonlinear prediction via Hermite transformation},
Statistical Theory and Related Fields 5 (2021), 49--54.
\href{https://doi.org/10.1080/24754269.2020.1856589}{doi:10.1080/24754269.2020.1856589}.
\bibitem{SV} T. Sottinen and L. Viitasaari,
\emph{Prediction law of fractional Brownian motion},
Statist. Probab. Lett. 129 (2017), 155--166.
\href{https://arxiv.org/abs/1610.02900}{arXiv:1610.02900}.
\bibitem{Yoshida} N. Yoshida, \emph{Asymptotic expansion of a variation with anticipative weights}, primary preprint, 2020, introduction and equations (1.2)--(1.7). \href{https://arxiv.org/abs/2101.00089}{arXiv:2101.00089}.
\bibitem{Pitera} M. Pitera,
\emph{Stochastic Processes---Lecture Notes}, Jagiellonian University,
November 29, 2020, Section 2.3, footnote 7, printed page 10.
\href{https://www2.im.uj.edu.pl/MarcinPitera/files/StochasticProcesses.pdf}{Primary author notes}.
\bibitem{Lachal} A. Lachal, \emph{A class of bridges of iterated
integrals of Brownian motion related to various boundary value problems
involving the one-dimensional polyharmonic operator}, International
Journal of Stochastic Analysis (2011), article 762486.
\href{https://math.univ-lyon1.fr/~alachal/english_publications.html}
{Primary author publication record and abstract}; detailed bridge
theorems not imported.
\bibitem{LP} D. A. Levin and Y. Peres, \emph{Markov Chains and Mixing Times}, second edition, with contributions by E. L. Wilmer, Section 13.2, Lemma 13.6. Author-hosted PDF, \href{https://pages.uoregon.edu/dlevin/MARKOV/markovmixing.pdf}{pages.uoregon.edu/dlevin/MARKOV/markovmixing.pdf}.
\end{thebibliography}
\end{document}
